\section{Scelta probabilistica e massa irrisolta}

Consideriamo le alternative alla prima decisione: quale consegna visitare
dopo il percorso iniziale. Ogni scelta identifica un sottoalbero di
completamenti. Trasformiamo i costi totali previsti in una distribuzione:
\[
 \pi(a\mid s_0)=\frac{\exp[-(\widehat f_a-m)/T]}
 {\sum_b\exp[-(\widehat f_b-m)/T]},
 \qquad m=\min_b\widehat f_b.
\]
La temperatura $T$ è la distanza media fra coppie di luoghi; vale uno
se tutte le distanze sono nulle. Queste probabilità esprimono preferenze
di esplorazione: costi stimati minori ricevono peso maggiore. Non sono
confidenze calibrate sull'ottimalità. La ricerca resta deterministica:
non campiona azioni. Cambiare $T$ modifica i pesi, ma non l'ordine delle stime.

Un'azione resta irrisolta se il suo sottoalbero contiene almeno uno
stato della frontiera attiva. La massa irrisolta è la somma delle
probabilità di queste azioni, ciascuna contata una volta:
\[
 \rho=\sum_{a\in U}\pi(a\mid s_0).
\]
Trovare un giro in un ramo non risolve tutti gli altri completamenti
possibili di quel ramo. La distribuzione iniziale rimane fissa durante
la prova; il cambiamento di massa dipende dalla ricerca. La stessa
misura è applicata ai tre sistemi, ma la baseline non la usa per scegliere
e il regolatore continua a osservare i due contatori. Una massa positiva
non è una probabilità di aver mancato l'ottimo e non equivale al numero
di stati pendenti.
